% !TeX root = ../../Book1.tex
% 纲要标题：### 2.6 射影映射与参数空间
% 纲要位置：几何学计划纲要.md，第 192 行。
% 正文状态：初稿；尚未正式审核。

% 纲要内容开始
% 
% * Segre 与二次 Veronese 映射的构造
% * Grassmann 空间与 Plücker 坐标
% * 点与超平面的关联簇及其投影
% 
% ---
% 
% 纲要内容结束

\section{射影映射与参数空间}
\label{g1:sec:0206}

圆锥曲线的参数化已经展示过一种变化：把一个点的坐标写成参数的二次式，便把一条射影线放进了另一个射影空间。还可以反过来问：一对射影点、一条直线或一个超平面，能否也用一个射影点来记录？合适的坐标会把“来自同一个对象”的相容条件写成方程，从而让整族几何对象本身成为可以研究的几何集合。本节始终取 \(K=\mathbb R\) 或 \(\mathbb C\)，只用多项式、矩阵秩与坐标图完成这些构造。

\subsection{Segre 与二次 Veronese 映射的构造}
\label{g1:subsec:0206:01}

先考虑两个独立的射影点 \([x_0:\cdots:x_m]\) 和 \([y_0:\cdots:y_n]\)。直接把坐标并排写成 \([x_0:\cdots:x_m:y_0:\cdots:y_n]\) 并不合法：改变第一组的代表而保持第二组不变，通常会改变所得射影点。若改用所有乘积 \(x_i y_j\)，两组代表分别乘以 \(\lambda\) 和 \(\mu\) 时，全部新坐标都乘以同一个 \(\lambda\mu\)。这正是齐次坐标允许的变化。

\begin{definition}[Segre 映射]\label{g1:def:segre-map}
	设 \(K=\mathbb R\) 或 \(\mathbb C\)，\(m,n\ge0\)。映射
	\[
		\sigma_{m,n}:\mathbb P^m(K)\times\mathbb P^n(K)
		\longrightarrow\mathbb P^{(m+1)(n+1)-1}(K),\qquad
		([x],[y])\longmapsto[x_i y_j]_{0\le i\le m,\,0\le j\le n}
	\]
	称为\term[Segreyingshe]{Segre 映射}[Segre map]，目标坐标先按 \(i\)、再按 \(j\) 排列。\footnote{Corrado Segre（1863--1924），意大利数学家，研究高维射影几何、代数曲面和直纹几何。}
\end{definition}

\begin{proposition}[乘积坐标与秩一矩阵]\label{g1:prop:segre-rank-one}
	设 \(K=\mathbb R\) 或 \(\mathbb C\)，\(m,n\ge0\)。Segre 映射良定义且为单射。把目标点的坐标排成 \((m+1)\times(n+1)\) 矩阵 \(Z=(z_{ij})\)，则其像恰由方程
	\[
		z_{ij}z_{k\ell}-z_{i\ell}z_{kj}=0
		\qquad(0\le i,k\le m,\ 0\le j,\ell\le n)
	\]
	给出，也就是所有非零秩一矩阵的射影类。在 \(z_{ab}\ne0\) 的部分，原来两个点可由第 \(b\) 列和第 \(a\) 行恢复。
\end{proposition}
\begin{proof}
	选取非零坐标 \(x_a,y_b\)，则 \(x_a y_b\ne0\)，所以输出坐标不全为零；代表的缩放只产生共同倍数，因而映射良定义。其矩阵是列向量与行向量的乘积 \(Z=xy^{\mathsf T}\)，每个二阶子式均为零。

	反过来，设 \(Z\ne0\) 且全部二阶子式为零，选定 \(z_{ab}\ne0\)。含第 \(a,i\) 行和第 \(b,j\) 列的子式给出
	\[
		z_{ij}=\frac{z_{ib}z_{aj}}{z_{ab}}.
	\]
	因此取 \(x_i=z_{ib}\)、\(y_j=z_{aj}/z_{ab}\)，便有 \(Z=xy^{\mathsf T}\)。这里 \(x_a\ne0\)、\(y_b=1\)，两向量均非零，证明了像的描述。对这样的分解，第 \(b\) 列与 \(x\) 成比例，第 \(a\) 行与 \(y^{\mathsf T}\) 成比例，故它们的射影类唯一，映射为单射。更具体地，在这张坐标图上，逆映射使用的仿射坐标为
	\[
		\frac{x_i}{x_a}=\frac{z_{ib}}{z_{ab}},\qquad
		\frac{y_j}{y_b}=\frac{z_{aj}}{z_{ab}}.
	\]
	分母在所选图上不为零，故这些公式也说明不同图中的恢复结果相容。
\end{proof}

\begin{example}[乘积中的两个方向]\label{g1:exm:segre-fibers}
	取 \(m=n=1\)，则
	\[
		([s:t],[u:v])\longmapsto[su:sv:tu:tv]
	\]
	将 \(\mathbb P^1(K)\times\mathbb P^1(K)\) 一一对应到 \(z_{00}z_{11}-z_{01}z_{10}=0\) 的二次曲面。固定 \([s:t]\) 而改变 \([u:v]\)，像是由 \((s,0,t,0)\)、\((0,s,0,t)\) 张成的射影直线；这两个向量线性无关。固定第二个因子同样得到一条直线。同一族中两条不同直线不相交，因为单射性不允许第一因子或第二因子同时取两个不同值；两族各取一条则恰交于相应有序对的像。\cref{g1:prop:quadric-two-rulings}中的两族直纹，在这里变成了乘积的两个独立方向。
\end{example}

还有另一种坐标构造。若希望把一个点的所有二次表达式同时作为坐标，就应取 \(x_i x_j\) 而不是取两组互不相关的变量。由此得到二次 Veronese 映射。\footnote{Giuseppe Veronese（1854--1917），意大利数学家，研究高维射影几何与几何基础。}

\begin{definition}[二次 Veronese 映射]\label{g1:def:veronese-map}
	设 \(K=\mathbb R\) 或 \(\mathbb C\)，\(n\ge1\)，\(N=\binom{n+2}{2}-1\)。映射
	\[
		\nu_2:\mathbb P^n(K)\longrightarrow\mathbb P^N(K),\qquad
		[x_0:\cdots:x_n]\longmapsto[x_i x_j]_{0\le i\le j\le n}
	\]
	称为\term[Veroneseyingshe]{二次 Veronese 映射}[quadratic Veronese map]。目标坐标按指标 \((i,j)\) 排列，将 \(z_{ij}\) 对称地补成矩阵 \(Z=(z_{ij})\)。
\end{definition}

\begin{proposition}[二次单项式坐标]\label{g1:prop:veronese-image}
	设 \(K=\mathbb R\) 或 \(\mathbb C\)，\(n\ge1\)。二次 Veronese 映射良定义且为单射，其像恰为非零对称秩一矩阵的射影类，亦即对称矩阵 \(Z\) 的所有二阶子式同时为零的射影点集。在 \(z_{aa}\ne0\) 的部分，逆映射由 \(x_i/x_a=z_{ia}/z_{aa}\) 给出。
\end{proposition}
\begin{proof}
	非零向量至少有一个非零坐标 \(x_a\)，故 \(x_a^2\ne0\)；缩放 \(x\) 使全部二次单项式乘以相同的平方倍数。这证明良定义，且像矩阵 \(xx^{\mathsf T}\) 为对称秩一矩阵。

	设对称矩阵 \(Z\ne0\) 的所有二阶子式均为零。必有某个对角元 \(z_{aa}\ne0\)：否则对任意 \(i,j\)，主子式 \(z_{ii}z_{jj}-z_{ij}^2=0\) 会强迫 \(z_{ij}=0\)，与 \(Z\ne0\) 矛盾。令 \(x_i=z_{ia}/z_{aa}\)，于是 \(x_a=1\)，而子式关系给出
	\[
		\frac{z_{ij}}{z_{aa}}
		=\frac{z_{ia}}{z_{aa}}\frac{z_{ja}}{z_{aa}}=x_i x_j.
	\]
	所以 \([Z]=[xx^{\mathsf T}]\)，并且上述比值唯一恢复 \([x]\)。这同时证明了像的描述、单射性及逆坐标公式。
\end{proof}

在实数情形，上一证明没有要求 \(z_{aa}\) 有实平方根。因为讨论的是矩阵的射影类，可以先把整个矩阵除以 \(z_{aa}\)；这也解释了为什么负半定的秩一矩阵仍表示像中的点。当 \(n=1\) 时，\(\nu_2([s:t])=[s^2:st:t^2]\)，正是前一节的圆锥曲线参数化。当 \(n=2\) 时，得到 \(\mathbb P^5(K)\) 中的 Veronese 曲面。

\ProseParagraphBreak

这两种映射还有一个直接用途：使原来的高次方程成为新空间的线性方程。例如任意二次齐次式 \(q(x)=\sum_{i\le j}c_{ij}x_i x_j\) 都满足
\[
	q(x)=0\quad\Longleftrightarrow\quad
	\sum_{i\le j}c_{ij}z_{ij}=0\quad\text{在 }[z]=\nu_2([x])\text{ 上成立}.
\]
右边是一个超平面与 Veronese 像的相交条件。同样，双线性关系 \(\sum a_{ij}x_i y_j=0\) 经 Segre 映射成为超平面截面。目标空间中的子式关系刻画了哪些射影点属于这些映射的像。这些模型的代数几何推广可参看\cite[第I章, 第2, Exercises节 2.12--2.15]{Hartshorne1977}。

\subsection{Grassmann 空间与 Plücker 坐标}
\label{g1:subsec:0206:02}

射影空间以一维向量子空间为点。若把“维数一”换成“维数 \(k\)”，同样得到一个值得研究的对象集合。困难在于：一个子空间有许多组基，把基向量的坐标直接抄下来会记录过多的选择。行列式恰能把换基的影响压缩成一个共同倍数。

\begin{definition}[Grassmann 空间]\label{g1:def:grassmannian}
	设 \(K=\mathbb R\) 或 \(\mathbb C\)，\(V\) 为 \(n\) 维 \(K\)-向量空间，\(1\le k\le n\)。所有 \(k\) 维向量子空间组成的集合
	\[
		\operatorname{Gr}(k,V)=\{\,W\subseteq V:W\text{ 为向量子空间且 }\dim_K W=k\,\}
	\]
	称为\term[Grassmannkongjian]{Grassmann 空间}[Grassmannian]，当 \(V=K^n\) 时也记为 \(\operatorname{Gr}(k,n;K)\)。\footnote{Hermann Günter Grassmann（1809--1877），德国数学家，发展了向量运算和外代数，并以坐标研究不同维数的子空间。}
	\symidx{\(\operatorname{Gr}(k,V)\)：Grassmann 空间}
\end{definition}

这里的 \(k,n\) 都是向量空间维数；\(W\) 对应 \(\mathbb P(V)\) 中的 \((k-1)\) 维射影子空间。特别地，\(\operatorname{Gr}(1,V)=\mathbb P(V)\)，而 \(\operatorname{Gr}(2,4;K)\) 参数化射影三维空间中的所有直线。先把后一情形具体写出。选取二维子空间的有序基 \(u,v\)，将它们写成矩阵的两行，并定义六个子式
\[
	p_{ij}=u_i v_j-u_j v_i\qquad(1\le i<j\le4).
\]
它们称为该直线的\term[Pluckerzuobiao]{Plücker 坐标}[Plücker coordinates]。\footnote{Julius Plücker（1801--1868），德国数学家和物理学家，研究解析几何、直线几何及光谱。}矩阵的秩为 \(2\)，故至少一个二阶子式非零。若将有序基替换为 \((u',v')^{\mathsf T}=S(u,v)^{\mathsf T}\)，其中 \(S\in\operatorname{GL}_2(K)\)，每个子式都乘以 \(\det S\)，所以射影点 \([p_{12}:p_{13}:p_{14}:p_{23}:p_{24}:p_{34}]\) 与基的选择无关。

\begin{theorem}[射影三维空间中直线的坐标]\label{g1:thm:plucker-gr24}
	设 \(K=\mathbb R\) 或 \(\mathbb C\)。Plücker 坐标给出 \(\operatorname{Gr}(2,4;K)\) 到 \(\mathbb P^5(K)\) 的单射，其像恰为二次曲面
	\begin{equation}\label{g1:eq:plucker-gr24}
		p_{12}p_{34}-p_{13}p_{24}+p_{14}p_{23}=0.
	\end{equation}
	在 \(p_{12}\ne0\) 的部分，每个子空间有唯一的行基矩阵
	\[
		\begin{pmatrix}
			1&0&a&b\\
			0&1&c&d
		\end{pmatrix},
	\]
	其中 \(a,b,c,d\in K\) 可任意取值，相应射影坐标为 \([1:c:d:-a:-b:ad-bc]\)。
\end{theorem}
\begin{proof}
	先说明关系式从何而来。将 \(p_{ij}=u_i v_j-u_j v_i\) 代入左边，展开得
	\[
		\begin{aligned}
			&(u_1v_2-u_2v_1)(u_3v_4-u_4v_3)\\
			&\quad-(u_1v_3-u_3v_1)(u_2v_4-u_4v_2)\\
			&\quad+(u_1v_4-u_4v_1)(u_2v_3-u_3v_2)=0;
		\end{aligned}
	\]
	每个含两个 \(u\) 坐标和两个 \(v\) 坐标的单项式出现两次且符号相反。这证明必要性。

	若 \(p_{12}\ne0\)，基矩阵的前两列组成可逆矩阵 \(C\)。将整矩阵左乘 \(C^{-1}\)，得到所述行基矩阵，而且前两坐标为 \((1,0)\)、\((0,1)\) 的两行在该子空间中唯一。逐项计算子式即得 \((1,c,d,-a,-b,ad-bc)\)。反过来，对满足关系式且 \(p_{12}\ne0\) 的射影点，将全部坐标除以 \(p_{12}\)，记归一化后的坐标为 \(\widetilde p_{ij}\)。取
	\[
		a=-\widetilde p_{23},\quad b=-\widetilde p_{24},\quad
		c=\widetilde p_{13},\quad d=\widetilde p_{14}.
	\]
	对应行基矩阵的前五个子式就是指定坐标；第六个由
	\[
		ad-bc=\widetilde p_{13}\widetilde p_{24}
		-\widetilde p_{14}\widetilde p_{23}=\widetilde p_{34}
	\]
	保证一致。这给出了存在性，也说明相同坐标只能恢复同一个子空间。

	任意非零六元组都有某个 \(p_{ij}\ne0\)。交换坐标轴可使这对指标成为 \(1,2\)。这里交换时同时使用 \(p_{ji}=-p_{ij}\)；关系式的左边在交换任意相邻两个指标时只改变符号，故它的零条件保持不变。于是上述构造适用于每一张图，并在图的重叠处由恢复出的子空间的唯一性相容，证明全局结论。
\end{proof}

\begin{example}[由六个数恢复一条直线]\label{g1:exm:plucker-reconstruction}
	在 \(K=\mathbb R\) 或 \(\mathbb C\) 上，六元组 \([1:2:3:4:5:-2]\) 满足 \(-2-10+12=0\)。按\cref{g1:thm:plucker-gr24}，其二维子空间有行基
	\[
		u=(1,0,-4,-5),\qquad v=(0,1,2,3),
	\]
	故对应直线为 \([s:t:-4s+2t:-5s+3t]\)，其中 \([s:t]\in\mathbb P^1(K)\)。若把最后一个坐标改成 \(0\)，关系式左边变成 \(2\)，所得射影点便不代表任何直线。所以六个齐次坐标之间的二次关系是实质约束。
\end{example}

一般的 \(\operatorname{Gr}(k,n;K)\) 也有同样可操作的坐标。将一组基写成秩为 \(k\) 的 \(k\times n\) 矩阵 \(M\)，选择一个非零的 \(k\) 阶子式。若它来自前 \(k\) 列，左乘这 \(k\) 列组成的矩阵的逆，即得唯一形式
\[
	M=[I_k\mid Z],\qquad Z\in\operatorname{Mat}_{k\times(n-k)}(K).
\]
唯一性的理由是：若 \(S[I_k\mid Z]=[I_k\mid Z']\)，比较前 \(k\) 列便有 \(S=I_k\)，继而 \(Z=Z'\)。任意 \(Z\) 都给出秩为 \(k\) 的矩阵。因此，这一部分中的子空间可用 \(k(n-k)\) 个自由参数描述。其他非零子式给出其他图；在重叠处，变换是先选出新的可逆子矩阵 \(C\)，再作 \(C^{-1}M\)。逆矩阵的余子式公式说明，新参数是旧参数的有理函数，分母就是被假设非零的 \(\det C\)。

\ProseParagraphBreak

所有 \(k\) 阶子式在换基时仍同乘 \(\det S\)，故又给出一个射影点，落在 \(\mathbb P^{\binom nk-1}(K)\) 中。这些坐标确实能恢复子空间：在选定子式归一化为 \(1\) 后， \([I_k\mid Z]\) 中的每个 \(Z_{ij}\) 都是用相应外部列替换第 \(i\) 个单位列所得子式，至多差一个由列排序产生的符号。因此所得到的射影映射为单射。一般维数时子式之间还有一组二次关系；式\eqref{g1:eq:plucker-gr24}给出了其中最小的非平凡模型。这里的参数数目由坐标图中的自由参数直接给出。

\subsection{点与超平面的关联簇及其投影}
\label{g1:subsec:0206:03}

对偶空间把一个超平面视为一个点；若同时保留原空间中的点，就能把“点属于超平面”本身作为一个几何对象。设 \(V=K^{n+1}\)、\(n\ge1\)。对偶空间 \(V^*\) 是 \(V\) 上的线性函数全体，非零函数 \(\alpha\) 的核给出超平面 \(\mathbb P(\ker\alpha)\)。定义
\[
	\mathcal I_n=
	\{\,([x],[\alpha])\in\mathbb P(V)\times\mathbb P(V^*):\alpha(x)=0\,\},
\]
称为点与超平面的\term[guanliancu]{关联簇}[incidence variety]。这里先将“簇”理解为这组方程给出的几何集合，关于代数簇的结构与一般性质将在第五卷中系统建立。在对偶基下，若 \(\alpha\) 的坐标为 \((a_0,\ldots,a_n)\)，其方程就是
\[
	\sum_{i=0}^n a_i x_i=0.
\]
分别缩放 \(a\) 和 \(x\) 不改变零条件，所以定义与代表选择无关。

\begin{proposition}[关联集合的两个投影]\label{g1:prop:incidence-fibers}
	设 \(K=\mathbb R\) 或 \(\mathbb C\)，\(V=K^{n+1}\)、\(n\ge1\)，并记 \(\mathcal I_n=\{([x],[\alpha])\in\mathbb P(V)\times\mathbb P(V^*):\alpha(x)=0\}\)。从 \(\mathcal I_n\) 到 \(\mathbb P(V)\)、\(\mathbb P(V^*)\) 的坐标投影都为满射。固定 \([x]\) 后，其上方的集合是 \(\mathbb P(\{\alpha:\alpha(x)=0\})\)，固定 \([\alpha]\) 后，其上方的集合是 \(\mathbb P(\ker\alpha)\)；二者都同构于 \(\mathbb P^{n-1}(K)\)。经 Segre 映射，\(\mathcal I_n\) 恰成为 Segre 像与超平面 \(\sum_i z_{ii}=0\) 的交集。
\end{proposition}
\begin{proof}
	固定非零向量 \(x\)，线性映射 \(V^*\rightarrow K\)、\(\alpha\mapsto\alpha(x)\) 是满射：将 \(x\) 扩为一组基，令线性函数在 \(x\) 上取值 \(1\)、在其余基向量上取值 \(0\) 即可。其核因而为 \(n\) 维，非零核向量的射影类就是所有经过 \([x]\) 的超平面。固定非零 \(\alpha\) 时，\(\alpha:V\rightarrow K\) 也满射，所以 \(\ker\alpha\) 同样为 \(n\) 维。由于 \(n\ge1\)，这些核都含非零向量，故两个投影都满射。最后，Segre 坐标取 \(z_{ij}=x_i a_j\) 时，关联方程恰为 \(\sum_i z_{ii}=0\)。\cref{g1:prop:segre-rank-one}给出的单射性和逆映射保证此描述没有额外或遗漏的点。
\end{proof}

分别记这两个投影为 \(\pi_1([x],[\alpha])=[x]\)、\(\pi_2([x],[\alpha])=[\alpha]\)。

\begin{figure}[htbp]
	\centering
	% !TeX root = ../Book1.tex
\begin{tikzcd}[column sep=large,row sep=large]
	& \mathcal I_n \arrow[dl, "\pi_1"'] \arrow[dr, "\pi_2"] & \\
	\mathbb P(V) && \mathbb P(V^*)
\end{tikzcd}

	\caption[点与超平面的关联]{点与超平面的关联。左边一个点的上方记录所有经过它的超平面；右边一个超平面的上方记录其全部点。}
	\label{g1:fig:incidence-projections}
\end{figure}

\cref{g1:fig:incidence-projections}把同一个条件的两种读法并排呈现。当 \(n=2\) 时，固定点得到过该点的直线束，固定直线得到这条直线上的点列；当 \(n=1\) 时，两边的纤维都是一个点，因为射影直线上的超平面本身就是点。这里映射的\term[xianwei]{纤维}[fiber]指一个指定点的逆像。

\ProseParagraphBreak

这些方程也解释了关联为什么是射影不变量。若 \(g\in\operatorname{GL}(V)\)，同时变换
\[
	[x]\longmapsto[gx],\qquad [\alpha]\longmapsto[\alpha\circ g^{-1}],
\]
则 \((\alpha\circ g^{-1})(gx)=\alpha(x)\)。因此点属于超平面的条件在改变射影坐标后保持不变。这些模型提出了三个共同问题：哪些条件描述的是一个子空间而非某组基，哪些参数变化只改变描述而不改变对象，哪些局部坐标能覆盖整个对象族。本节的秩一分解、子式恢复和关联方程已经分别给出了可直接计算的答案。后续研究更一般的曲线、曲面和参数空间时，这三个问题仍然存在。
